\documentclass[11pt]{article}
\usepackage[margin=1.5cm]{geometry}
\usepackage{amsmath,amssymb}
\usepackage{enumitem}
\usepackage{tikz}
\usepackage{listings}
\usepackage[T1]{fontenc}
\setlength{\columnsep}{0.7cm}
\setlength{\parindent}{0pt}
\setlist{nosep,leftmargin=*}
\lstset{
  basicstyle=\ttfamily\scriptsize,
  columns=fullflexible,
  keepspaces=true,
  aboveskip=2pt,
  belowskip=2pt,
  frame=single,
  breaklines=true
}

\begin{document}
\twocolumn

\title{Chapter 13 In-Class Exercise: ComplexNumber Inheritance}
\author{Prof. Jordan C. Hanson}
\date{\today}
\maketitle

\section{Parent Class}

Assume \texttt{vector2d.py} contains:

\begin{verbatim}
import math
class Vector2D:
    def __init__(self, x, y):
        self.x = x
        self.y = y
    def magnitude(self):
        return math.sqrt(self.x**2+self.y**2)
    def angle(self):
        return math.atan2(self.y,self.x)
    def __add__(self, other):
        return Vector2D(
            self.x+other.x,self.y+other.y)
    def __sub__(self, other):
        return Vector2D(
            self.x-other.x,self.y-other.y)
\end{verbatim}

A complex number \(z=a+jb\) also has two components. We will use $x\leftrightarrow a,\qquad y\leftrightarrow b$.

\section{Create the Child Class}

Create \texttt{complex\_number.py}:

\begin{verbatim}
from vector2d import Vector2D
class ComplexNumber(Vector2D):
    def __init__(self, real, imag):
        super().__init__(____,____)
\end{verbatim}

The expression \texttt{ComplexNumber(Vector2D)} means that
\texttt{ComplexNumber} \textbf{inherits} from \texttt{Vector2D}. Test:

\begin{verbatim}
z1 = ComplexNumber(3, 4)
print(z1.magnitude())
print(z1.angle())
\end{verbatim}

\section{Override Addition}

Now try

\begin{verbatim}
z1 = ComplexNumber(3,2)
z2 = ComplexNumber(5,-4)
z3 = z1 + z2
print(type(z3))
\end{verbatim}

The arithmetic is correct, but the inherited method returns a \texttt{Vector2D}. Override it:

\begin{verbatim}
def __add__(self, other):
    return ComplexNumber(
        self.x + other.x,
        self.y + other.y)
\end{verbatim}

What should \texttt{type(z3)} give, if you recalculate $z1+z2$?  Override the subtraction method in the same way.

\section{New Behavior: Multiplication}

Complex multiplication is
\[
(a+jb)(c+jd)=(ac-bd)+(ad+bc)j.
\]

Add:

\begin{verbatim}
def __mul__(self, other):
    real = ____
    imag = ____
    return ComplexNumber(real,imag)
\end{verbatim}

Test the new method.

\section{Complex-Number Names}

The parent stores \texttt{x} and \texttt{y}. Add:

\begin{verbatim}
def real(self):
    return self.x
def imag(self):
    return self.y
\end{verbatim}

Then test

\begin{verbatim}
print(z1.real())
print(z1.imag())
\end{verbatim}

\textbf{Key idea:} Inheritance reuses behavior; overriding replaces
inherited behavior when the child class needs something different.

\end{document}
